% Archived from report/main.tex on 2025-06-14.
% Section removed from the report to reduce scope; restore with % Archived from report/main.tex on 2025-06-14.
% Section removed from the report to reduce scope; restore with % Archived from report/main.tex on 2025-06-14.
% Section removed from the report to reduce scope; restore with % Archived from report/main.tex on 2025-06-14.
% Section removed from the report to reduce scope; restore with \input{archive/multi-agent-frameworks}
% if needed.

\section{Multi-Agent Frameworks for Code Generation}
\label{sec:multiagent}

An alternative architectural approach to code generation distributes the work
across multiple specialised agents, each responsible for a distinct phase or
concern. Where single-agent systems integrate planning, generation, testing,
and refinement into one LLM conversation, multi-agent frameworks decompose the
software lifecycle into roles---architect, programmer, tester, reviewer---and
orchestrate collaboration among agents playing those roles. This
division-of-labour mirrors human software teams and aims to improve code
quality, modularity, and error detection beyond what monolithic agents achieve.

This section surveys role decomposition and communication patterns
(\S\ref{sec:multi-roles}), examines representative framework designs
(\S\ref{sec:multi-frameworks}), and assesses when multi-agent designs
outperform monolithic alternatives (\S\ref{sec:multi-when}).

\subsection{Role Decomposition and Communication}
\label{sec:multi-roles}

Multi-agent systems partition code-generation tasks along two primary
dimensions: \emph{temporal stages} (design, implementation, testing, review)
and \emph{functional concerns} (requirements analysis, algorithm design,
execution, debugging). Agents are assigned roles corresponding to these
partitions, and collaboration protocols define how information flows among them.

\paragraph{Role taxonomy.}
Common roles mirror human software teams. A \emph{product manager} or
\emph{architect} agent refines high-level requirements into technical
specifications. A \emph{programmer} or \emph{coder} agent generates
implementation. A \emph{test designer} agent writes test cases independently of
the code. A \emph{test executor} agent runs code against tests and reports
failures. A \emph{reviewer} agent critiques code for correctness, style, or
security. More specialised roles include a \emph{retrieval agent} that searches
documentation or past solutions, and a \emph{debugger agent} that localises and
repairs faults. Not all frameworks use all roles; designs range from three
agents to dozens.

\paragraph{Communication patterns.}
Agents exchange information through structured protocols. The simplest is
\emph{sequential handoff}: the architect produces a design document, passes it
to the programmer, who generates code, passes it to the tester, and so on. More
sophisticated patterns include \emph{iterative feedback loops}, where the test
executor returns failure traces to the programmer for revision, and
\emph{debate or voting}, where multiple agents propose alternatives and a
consensus mechanism selects the best. Communication may use natural language
(agents converse as humans would), structured artifacts (JSON schemas, API
specifications, test suites), or hybrid approaches.

\paragraph{Memory and state management.}
Multi-agent systems maintain shared state through blackboard architectures,
message queues, or centralised repositories. A blackboard allows all agents to
read and write a common workspace containing the current design, code, tests,
and execution logs. Message queues enforce ordering and ensure agents see
consistent snapshots of the task. Centralised repositories (e.g., a simulated
Git workflow) provide versioning and rollback. The choice of state management
affects coordination overhead: blackboards are flexible but require
synchronisation; message queues are deterministic but rigid; repositories add
version-control semantics but increase complexity.

\subsection{Framework Designs}
\label{sec:multi-frameworks}

We survey four representative multi-agent frameworks that span diverse design
choices in orchestration, role granularity, and evaluation focus.
Table~\ref{tab:multiagent} provides a comparative overview.

\begin{table*}[t]
  \centering
  \caption{Comparison of multi-agent code-generation frameworks.}
  \label{tab:multiagent}
  \small
  \begin{tabular}{@{}llp{0.28\textwidth}p{0.25\textwidth}cc@{}}
    \toprule
    \textbf{System} & \textbf{Roles} & \textbf{Communication} & \textbf{Key mechanism} & \textbf{HumanEval} & \textbf{MBPP} \\
    \midrule
    ChatDev & CEO, CPO, CTO, Programmer, Reviewer, Tester
      & Sequential chat-based handoffs with communicative dehallucination
      & Waterfall-inspired phases; memory replay to prevent hallucinations
      & -- & -- \\
    MetaGPT & Product manager, Architect, Engineer, QA
      & SOP-encoded prompts with structured artifacts (PRD, design docs)
      & Assembly-line paradigm; executable feedback module
      & 85.9\% & 87.7\% \\
    AgentCoder & Programmer, Test designer, Test executor
      & Iterative loop: code $\to$ tests $\to$ execution $\to$ feedback
      & Independent test generation; execution-driven refinement
      & 77.4\% & 89.1\% \\
    MapCoder & Retrieval, Planner, Coder, Debugger
      & Sequential pipeline with recalling $\to$ planning $\to$ coding $\to$ debugging
      & Mimics human problem-solving cycle; example retrieval
      & 93.9\% & 83.1\% \\
    \bottomrule
  \end{tabular}
\end{table*}

\paragraph{ChatDev~\cite{qian2024chatdev}.}
From Tsinghua University and collaborators, ChatDev models a virtual software
company in which agents assume executive and technical roles: CEO, Chief
Product Officer (CPO), Chief Technology Officer (CTO), Programmer, Reviewer,
Tester, and Art Designer. The system organises work into a waterfall-like
sequence of ``chat chains'': each phase (designing, coding, testing,
documenting) corresponds to a multi-turn conversation among relevant agents.
For example, the CPO and CEO refine requirements, then the CTO and Programmer
generate code, then the Programmer and Tester validate it. ChatDev introduces
\emph{communicative dehallucination}: agents maintain memory of earlier
discussions and replay relevant context to reduce inconsistencies and
hallucinations across phases. The framework also supports role-playing prompts
that instruct each LLM instantiation to adopt a persona, though all agents
share the same base model. ChatDev demonstrates end-to-end software generation
from natural-language descriptions to executable programs, but performance
metrics on standardised benchmarks are limited in the original publication.

\paragraph{MetaGPT~\cite{hong2024metagpt}.}
MetaGPT, from DeepWisdom and collaborators, encodes human-like Standardized
Operating Procedures (SOPs) into LLM prompts to streamline multi-agent
collaboration. It assigns four primary roles: Product Manager, Architect,
Engineer, and QA. Each role produces structured artifacts---Product Requirement
Documents (PRD), system designs, implementation code, and test plans---passed
sequentially along an ``assembly line''. The key innovation is that agents
communicate via formalised documents rather than free-form chat, reducing
ambiguity and cascading hallucinations. MetaGPT includes an \emph{executable
feedback module}: generated code is run against tests, and failures trigger
targeted revisions by the Engineer agent. On HumanEval and MBPP, MetaGPT
achieves 85.9\% and 87.7\% pass@1 respectively, outperforming many single-agent
baselines and prior multi-agent systems. The framework demonstrates that
structured inter-agent protocols can significantly reduce errors compared to
informal chat-based collaboration.

\paragraph{AgentCoder~\cite{huang2023agentcoder}.}
AgentCoder, from the University of Hong Kong and collaborators, simplifies role
structure to three agents: Programmer, Test Designer, and Test Executor. The
Programmer generates code. The Test Designer independently generates test
cases, ensuring tests are not biased by the code. The Test Executor runs code
against tests and reports failures. If tests fail, the executor provides
feedback (error messages, stack traces) to the Programmer, which revises the
code. This cycle iterates until all tests pass or a budget is exhausted.
AgentCoder's strength lies in \emph{independent test generation}: because tests
are written without seeing the code, they better reflect specification
requirements rather than implementation quirks. On HumanEval and MBPP,
AgentCoder with GPT-3.5 reaches 77.4\% and 89.1\% pass@1, surpassing
single-agent GPT-3.5 baselines (69.5\%, 63.0\%) and demonstrating that even a
minimal three-agent decomposition improves reliability.

\paragraph{MapCoder~\cite{islam2024mapcoder}.}
MapCoder, from Bangladesh University of Engineering and UCLA, targets
competitive programming by mimicking the human problem-solving cycle through
four agents: Recall, Plan, Code, and Debug. The Recall agent retrieves similar
past problems or solution templates from a knowledge base. The Plan agent
decomposes the problem into algorithmic steps. The Code agent generates the
implementation. The Debug agent runs the code, identifies failures, and
proposes fixes. Agents operate sequentially, passing intermediate results
forward, and backtracking occurs when debugging reveals a flaw in earlier
stages. MapCoder achieves strong results on competitive benchmarks: 93.9\%
pass@1 on HumanEval, 83.1\% on MBPP, and 28.5\% on CodeContests. The framework
shows that task-specific role design (e.g., retrieval for competitive
programming) can yield state-of-the-art performance when aligned with domain
workflows.

\subsection{When Does Multi-Agency Help?}
\label{sec:multi-when}

Multi-agent designs introduce coordination overhead, increased token costs, and
additional failure modes. Whether they outperform monolithic agents depends on
task complexity, the quality of role decomposition, and the effectiveness of
inter-agent communication.

\paragraph{Benefits: specialisation and modularity.}
The primary advantage of multi-agent systems is \emph{specialisation}. When a
single agent must simultaneously plan, code, test, and debug, its attention is
divided and errors in one phase can cascade unchecked. Decomposing these
concerns into separate agents allows each to focus on a narrower task. For
example, AgentCoder's independent test designer generates tests that are more
comprehensive and less biased than tests written by the same agent that wrote
the code. Similarly, MetaGPT's structured artifacts (PRDs, design documents)
force explicit planning before coding, reducing the likelihood that
implementation drifts from requirements. Modularity also enables targeted
improvement: if test generation is weak, only the test-designer agent need be
upgraded, rather than retraining or reprompting the entire system.

\paragraph{Benefits: error detection and correction.}
Multi-agent systems create multiple checkpoints where errors can be caught. A
reviewer agent can identify code smells or logic bugs before execution. A
separate test executor can detect runtime failures and provide diagnostic
feedback to the programmer without the programmer needing to self-critique.
This separation mirrors human code review and pair programming, which are
empirically known to improve code quality. MapCoder's debug agent, for example,
systematically analyses test failures and proposes fixes, iterating until the
code passes. Single-agent systems can approximate this through
self-reflection~\cite{madaan2023selfrefine}, but multi-agent designs make the
separation explicit and avoid the risk of self-confirming biases.

\paragraph{Costs: coordination and token overhead.}
The main cost is coordination complexity. Each inter-agent communication
requires prompt construction, LLM inference, and state synchronisation. If
agents operate sequentially, total latency is the sum of all agent runtimes. If
agents operate in parallel (e.g., multiple programmers proposing solutions),
orchestration logic must merge or vote on outputs. Token costs scale with the
number of agents and communication rounds: MetaGPT and ChatDev can consume
several times the tokens of a single-agent baseline. In practice, frameworks
mitigate this by caching shared context, compressing intermediate artifacts, or
limiting the number of iterations. However, for simple tasks where a
single-agent solution suffices, the overhead is not justified.

\paragraph{Failure modes: handoff errors and misalignment.}
Multi-agent systems introduce new failure modes. Handoff errors occur when one
agent's output is misunderstood or ignored by the next. For example, if the
architect agent specifies an API contract in natural language and the
programmer misinterprets it, the resulting code may be internally consistent
but wrong. MetaGPT's structured artifacts reduce this risk by enforcing
explicit schemas, but free-form chat-based systems like ChatDev remain
vulnerable. Another failure mode is \emph{agent misalignment}: if the test
designer agent generates trivial or incorrect tests, the system may converge on
code that passes those weak tests but fails in production. AgentCoder addresses
this by ensuring tests are written independently, but the quality still depends
on the test-designer prompt and model.

\paragraph{Diminishing returns with agent count.}
Adding more agents does not monotonically improve performance. Beyond a certain
point, coordination overhead dominates and additional roles contribute little.
Empirical studies on ChatDev and MetaGPT suggest that 3--5 agents (planning,
coding, testing, review) capture most of the benefit, while systems with 10+
agents show marginal improvements at much higher cost. The optimal agent count
depends on task complexity: simple function-writing benefits from
programmer-tester pairs, while full software projects may justify
architect-programmer-tester-reviewer-documenter chains. However, even for
complex tasks, the evidence does not yet support very large agent counts
analogous to human teams of dozens.

\paragraph{When to choose multi-agent designs.}
Multi-agent frameworks are most beneficial when: (i)~the task involves multiple
distinct phases (design, implementation, validation) that benefit from
specialisation; (ii)~intermediate artifacts (specifications, test suites) can
be produced and validated independently; (iii)~the base LLM is strong enough
that role-specific prompting yields measurably different behaviours; and
(iv)~the cost overhead is acceptable given the quality improvement. For
simple one-off code snippets, single-agent systems with self-refinement suffice
and are faster. For end-to-end software engineering tasks on benchmarks like
SWE-bench, the evidence is mixed: some multi-agent systems underperform
well-tuned single-agent scaffolds (recall Agentless outperforming many agentic
designs in \S\ref{sec:e2e}), suggesting that scaffold quality matters as much
as agent count. The frontier remains open, and multi-agent research continues
to refine orchestration protocols and role definitions.

% if needed.

\section{Multi-Agent Frameworks for Code Generation}
\label{sec:multiagent}

An alternative architectural approach to code generation distributes the work
across multiple specialised agents, each responsible for a distinct phase or
concern. Where single-agent systems integrate planning, generation, testing,
and refinement into one LLM conversation, multi-agent frameworks decompose the
software lifecycle into roles---architect, programmer, tester, reviewer---and
orchestrate collaboration among agents playing those roles. This
division-of-labour mirrors human software teams and aims to improve code
quality, modularity, and error detection beyond what monolithic agents achieve.

This section surveys role decomposition and communication patterns
(\S\ref{sec:multi-roles}), examines representative framework designs
(\S\ref{sec:multi-frameworks}), and assesses when multi-agent designs
outperform monolithic alternatives (\S\ref{sec:multi-when}).

\subsection{Role Decomposition and Communication}
\label{sec:multi-roles}

Multi-agent systems partition code-generation tasks along two primary
dimensions: \emph{temporal stages} (design, implementation, testing, review)
and \emph{functional concerns} (requirements analysis, algorithm design,
execution, debugging). Agents are assigned roles corresponding to these
partitions, and collaboration protocols define how information flows among them.

\paragraph{Role taxonomy.}
Common roles mirror human software teams. A \emph{product manager} or
\emph{architect} agent refines high-level requirements into technical
specifications. A \emph{programmer} or \emph{coder} agent generates
implementation. A \emph{test designer} agent writes test cases independently of
the code. A \emph{test executor} agent runs code against tests and reports
failures. A \emph{reviewer} agent critiques code for correctness, style, or
security. More specialised roles include a \emph{retrieval agent} that searches
documentation or past solutions, and a \emph{debugger agent} that localises and
repairs faults. Not all frameworks use all roles; designs range from three
agents to dozens.

\paragraph{Communication patterns.}
Agents exchange information through structured protocols. The simplest is
\emph{sequential handoff}: the architect produces a design document, passes it
to the programmer, who generates code, passes it to the tester, and so on. More
sophisticated patterns include \emph{iterative feedback loops}, where the test
executor returns failure traces to the programmer for revision, and
\emph{debate or voting}, where multiple agents propose alternatives and a
consensus mechanism selects the best. Communication may use natural language
(agents converse as humans would), structured artifacts (JSON schemas, API
specifications, test suites), or hybrid approaches.

\paragraph{Memory and state management.}
Multi-agent systems maintain shared state through blackboard architectures,
message queues, or centralised repositories. A blackboard allows all agents to
read and write a common workspace containing the current design, code, tests,
and execution logs. Message queues enforce ordering and ensure agents see
consistent snapshots of the task. Centralised repositories (e.g., a simulated
Git workflow) provide versioning and rollback. The choice of state management
affects coordination overhead: blackboards are flexible but require
synchronisation; message queues are deterministic but rigid; repositories add
version-control semantics but increase complexity.

\subsection{Framework Designs}
\label{sec:multi-frameworks}

We survey four representative multi-agent frameworks that span diverse design
choices in orchestration, role granularity, and evaluation focus.
Table~\ref{tab:multiagent} provides a comparative overview.

\begin{table*}[t]
  \centering
  \caption{Comparison of multi-agent code-generation frameworks.}
  \label{tab:multiagent}
  \small
  \begin{tabular}{@{}llp{0.28\textwidth}p{0.25\textwidth}cc@{}}
    \toprule
    \textbf{System} & \textbf{Roles} & \textbf{Communication} & \textbf{Key mechanism} & \textbf{HumanEval} & \textbf{MBPP} \\
    \midrule
    ChatDev & CEO, CPO, CTO, Programmer, Reviewer, Tester
      & Sequential chat-based handoffs with communicative dehallucination
      & Waterfall-inspired phases; memory replay to prevent hallucinations
      & -- & -- \\
    MetaGPT & Product manager, Architect, Engineer, QA
      & SOP-encoded prompts with structured artifacts (PRD, design docs)
      & Assembly-line paradigm; executable feedback module
      & 85.9\% & 87.7\% \\
    AgentCoder & Programmer, Test designer, Test executor
      & Iterative loop: code $\to$ tests $\to$ execution $\to$ feedback
      & Independent test generation; execution-driven refinement
      & 77.4\% & 89.1\% \\
    MapCoder & Retrieval, Planner, Coder, Debugger
      & Sequential pipeline with recalling $\to$ planning $\to$ coding $\to$ debugging
      & Mimics human problem-solving cycle; example retrieval
      & 93.9\% & 83.1\% \\
    \bottomrule
  \end{tabular}
\end{table*}

\paragraph{ChatDev~\cite{qian2024chatdev}.}
From Tsinghua University and collaborators, ChatDev models a virtual software
company in which agents assume executive and technical roles: CEO, Chief
Product Officer (CPO), Chief Technology Officer (CTO), Programmer, Reviewer,
Tester, and Art Designer. The system organises work into a waterfall-like
sequence of ``chat chains'': each phase (designing, coding, testing,
documenting) corresponds to a multi-turn conversation among relevant agents.
For example, the CPO and CEO refine requirements, then the CTO and Programmer
generate code, then the Programmer and Tester validate it. ChatDev introduces
\emph{communicative dehallucination}: agents maintain memory of earlier
discussions and replay relevant context to reduce inconsistencies and
hallucinations across phases. The framework also supports role-playing prompts
that instruct each LLM instantiation to adopt a persona, though all agents
share the same base model. ChatDev demonstrates end-to-end software generation
from natural-language descriptions to executable programs, but performance
metrics on standardised benchmarks are limited in the original publication.

\paragraph{MetaGPT~\cite{hong2024metagpt}.}
MetaGPT, from DeepWisdom and collaborators, encodes human-like Standardized
Operating Procedures (SOPs) into LLM prompts to streamline multi-agent
collaboration. It assigns four primary roles: Product Manager, Architect,
Engineer, and QA. Each role produces structured artifacts---Product Requirement
Documents (PRD), system designs, implementation code, and test plans---passed
sequentially along an ``assembly line''. The key innovation is that agents
communicate via formalised documents rather than free-form chat, reducing
ambiguity and cascading hallucinations. MetaGPT includes an \emph{executable
feedback module}: generated code is run against tests, and failures trigger
targeted revisions by the Engineer agent. On HumanEval and MBPP, MetaGPT
achieves 85.9\% and 87.7\% pass@1 respectively, outperforming many single-agent
baselines and prior multi-agent systems. The framework demonstrates that
structured inter-agent protocols can significantly reduce errors compared to
informal chat-based collaboration.

\paragraph{AgentCoder~\cite{huang2023agentcoder}.}
AgentCoder, from the University of Hong Kong and collaborators, simplifies role
structure to three agents: Programmer, Test Designer, and Test Executor. The
Programmer generates code. The Test Designer independently generates test
cases, ensuring tests are not biased by the code. The Test Executor runs code
against tests and reports failures. If tests fail, the executor provides
feedback (error messages, stack traces) to the Programmer, which revises the
code. This cycle iterates until all tests pass or a budget is exhausted.
AgentCoder's strength lies in \emph{independent test generation}: because tests
are written without seeing the code, they better reflect specification
requirements rather than implementation quirks. On HumanEval and MBPP,
AgentCoder with GPT-3.5 reaches 77.4\% and 89.1\% pass@1, surpassing
single-agent GPT-3.5 baselines (69.5\%, 63.0\%) and demonstrating that even a
minimal three-agent decomposition improves reliability.

\paragraph{MapCoder~\cite{islam2024mapcoder}.}
MapCoder, from Bangladesh University of Engineering and UCLA, targets
competitive programming by mimicking the human problem-solving cycle through
four agents: Recall, Plan, Code, and Debug. The Recall agent retrieves similar
past problems or solution templates from a knowledge base. The Plan agent
decomposes the problem into algorithmic steps. The Code agent generates the
implementation. The Debug agent runs the code, identifies failures, and
proposes fixes. Agents operate sequentially, passing intermediate results
forward, and backtracking occurs when debugging reveals a flaw in earlier
stages. MapCoder achieves strong results on competitive benchmarks: 93.9\%
pass@1 on HumanEval, 83.1\% on MBPP, and 28.5\% on CodeContests. The framework
shows that task-specific role design (e.g., retrieval for competitive
programming) can yield state-of-the-art performance when aligned with domain
workflows.

\subsection{When Does Multi-Agency Help?}
\label{sec:multi-when}

Multi-agent designs introduce coordination overhead, increased token costs, and
additional failure modes. Whether they outperform monolithic agents depends on
task complexity, the quality of role decomposition, and the effectiveness of
inter-agent communication.

\paragraph{Benefits: specialisation and modularity.}
The primary advantage of multi-agent systems is \emph{specialisation}. When a
single agent must simultaneously plan, code, test, and debug, its attention is
divided and errors in one phase can cascade unchecked. Decomposing these
concerns into separate agents allows each to focus on a narrower task. For
example, AgentCoder's independent test designer generates tests that are more
comprehensive and less biased than tests written by the same agent that wrote
the code. Similarly, MetaGPT's structured artifacts (PRDs, design documents)
force explicit planning before coding, reducing the likelihood that
implementation drifts from requirements. Modularity also enables targeted
improvement: if test generation is weak, only the test-designer agent need be
upgraded, rather than retraining or reprompting the entire system.

\paragraph{Benefits: error detection and correction.}
Multi-agent systems create multiple checkpoints where errors can be caught. A
reviewer agent can identify code smells or logic bugs before execution. A
separate test executor can detect runtime failures and provide diagnostic
feedback to the programmer without the programmer needing to self-critique.
This separation mirrors human code review and pair programming, which are
empirically known to improve code quality. MapCoder's debug agent, for example,
systematically analyses test failures and proposes fixes, iterating until the
code passes. Single-agent systems can approximate this through
self-reflection~\cite{madaan2023selfrefine}, but multi-agent designs make the
separation explicit and avoid the risk of self-confirming biases.

\paragraph{Costs: coordination and token overhead.}
The main cost is coordination complexity. Each inter-agent communication
requires prompt construction, LLM inference, and state synchronisation. If
agents operate sequentially, total latency is the sum of all agent runtimes. If
agents operate in parallel (e.g., multiple programmers proposing solutions),
orchestration logic must merge or vote on outputs. Token costs scale with the
number of agents and communication rounds: MetaGPT and ChatDev can consume
several times the tokens of a single-agent baseline. In practice, frameworks
mitigate this by caching shared context, compressing intermediate artifacts, or
limiting the number of iterations. However, for simple tasks where a
single-agent solution suffices, the overhead is not justified.

\paragraph{Failure modes: handoff errors and misalignment.}
Multi-agent systems introduce new failure modes. Handoff errors occur when one
agent's output is misunderstood or ignored by the next. For example, if the
architect agent specifies an API contract in natural language and the
programmer misinterprets it, the resulting code may be internally consistent
but wrong. MetaGPT's structured artifacts reduce this risk by enforcing
explicit schemas, but free-form chat-based systems like ChatDev remain
vulnerable. Another failure mode is \emph{agent misalignment}: if the test
designer agent generates trivial or incorrect tests, the system may converge on
code that passes those weak tests but fails in production. AgentCoder addresses
this by ensuring tests are written independently, but the quality still depends
on the test-designer prompt and model.

\paragraph{Diminishing returns with agent count.}
Adding more agents does not monotonically improve performance. Beyond a certain
point, coordination overhead dominates and additional roles contribute little.
Empirical studies on ChatDev and MetaGPT suggest that 3--5 agents (planning,
coding, testing, review) capture most of the benefit, while systems with 10+
agents show marginal improvements at much higher cost. The optimal agent count
depends on task complexity: simple function-writing benefits from
programmer-tester pairs, while full software projects may justify
architect-programmer-tester-reviewer-documenter chains. However, even for
complex tasks, the evidence does not yet support very large agent counts
analogous to human teams of dozens.

\paragraph{When to choose multi-agent designs.}
Multi-agent frameworks are most beneficial when: (i)~the task involves multiple
distinct phases (design, implementation, validation) that benefit from
specialisation; (ii)~intermediate artifacts (specifications, test suites) can
be produced and validated independently; (iii)~the base LLM is strong enough
that role-specific prompting yields measurably different behaviours; and
(iv)~the cost overhead is acceptable given the quality improvement. For
simple one-off code snippets, single-agent systems with self-refinement suffice
and are faster. For end-to-end software engineering tasks on benchmarks like
SWE-bench, the evidence is mixed: some multi-agent systems underperform
well-tuned single-agent scaffolds (recall Agentless outperforming many agentic
designs in \S\ref{sec:e2e}), suggesting that scaffold quality matters as much
as agent count. The frontier remains open, and multi-agent research continues
to refine orchestration protocols and role definitions.

% if needed.

\section{Multi-Agent Frameworks for Code Generation}
\label{sec:multiagent}

An alternative architectural approach to code generation distributes the work
across multiple specialised agents, each responsible for a distinct phase or
concern. Where single-agent systems integrate planning, generation, testing,
and refinement into one LLM conversation, multi-agent frameworks decompose the
software lifecycle into roles---architect, programmer, tester, reviewer---and
orchestrate collaboration among agents playing those roles. This
division-of-labour mirrors human software teams and aims to improve code
quality, modularity, and error detection beyond what monolithic agents achieve.

This section surveys role decomposition and communication patterns
(\S\ref{sec:multi-roles}), examines representative framework designs
(\S\ref{sec:multi-frameworks}), and assesses when multi-agent designs
outperform monolithic alternatives (\S\ref{sec:multi-when}).

\subsection{Role Decomposition and Communication}
\label{sec:multi-roles}

Multi-agent systems partition code-generation tasks along two primary
dimensions: \emph{temporal stages} (design, implementation, testing, review)
and \emph{functional concerns} (requirements analysis, algorithm design,
execution, debugging). Agents are assigned roles corresponding to these
partitions, and collaboration protocols define how information flows among them.

\paragraph{Role taxonomy.}
Common roles mirror human software teams. A \emph{product manager} or
\emph{architect} agent refines high-level requirements into technical
specifications. A \emph{programmer} or \emph{coder} agent generates
implementation. A \emph{test designer} agent writes test cases independently of
the code. A \emph{test executor} agent runs code against tests and reports
failures. A \emph{reviewer} agent critiques code for correctness, style, or
security. More specialised roles include a \emph{retrieval agent} that searches
documentation or past solutions, and a \emph{debugger agent} that localises and
repairs faults. Not all frameworks use all roles; designs range from three
agents to dozens.

\paragraph{Communication patterns.}
Agents exchange information through structured protocols. The simplest is
\emph{sequential handoff}: the architect produces a design document, passes it
to the programmer, who generates code, passes it to the tester, and so on. More
sophisticated patterns include \emph{iterative feedback loops}, where the test
executor returns failure traces to the programmer for revision, and
\emph{debate or voting}, where multiple agents propose alternatives and a
consensus mechanism selects the best. Communication may use natural language
(agents converse as humans would), structured artifacts (JSON schemas, API
specifications, test suites), or hybrid approaches.

\paragraph{Memory and state management.}
Multi-agent systems maintain shared state through blackboard architectures,
message queues, or centralised repositories. A blackboard allows all agents to
read and write a common workspace containing the current design, code, tests,
and execution logs. Message queues enforce ordering and ensure agents see
consistent snapshots of the task. Centralised repositories (e.g., a simulated
Git workflow) provide versioning and rollback. The choice of state management
affects coordination overhead: blackboards are flexible but require
synchronisation; message queues are deterministic but rigid; repositories add
version-control semantics but increase complexity.

\subsection{Framework Designs}
\label{sec:multi-frameworks}

We survey four representative multi-agent frameworks that span diverse design
choices in orchestration, role granularity, and evaluation focus.
Table~\ref{tab:multiagent} provides a comparative overview.

\begin{table*}[t]
  \centering
  \caption{Comparison of multi-agent code-generation frameworks.}
  \label{tab:multiagent}
  \small
  \begin{tabular}{@{}llp{0.28\textwidth}p{0.25\textwidth}cc@{}}
    \toprule
    \textbf{System} & \textbf{Roles} & \textbf{Communication} & \textbf{Key mechanism} & \textbf{HumanEval} & \textbf{MBPP} \\
    \midrule
    ChatDev & CEO, CPO, CTO, Programmer, Reviewer, Tester
      & Sequential chat-based handoffs with communicative dehallucination
      & Waterfall-inspired phases; memory replay to prevent hallucinations
      & -- & -- \\
    MetaGPT & Product manager, Architect, Engineer, QA
      & SOP-encoded prompts with structured artifacts (PRD, design docs)
      & Assembly-line paradigm; executable feedback module
      & 85.9\% & 87.7\% \\
    AgentCoder & Programmer, Test designer, Test executor
      & Iterative loop: code $\to$ tests $\to$ execution $\to$ feedback
      & Independent test generation; execution-driven refinement
      & 77.4\% & 89.1\% \\
    MapCoder & Retrieval, Planner, Coder, Debugger
      & Sequential pipeline with recalling $\to$ planning $\to$ coding $\to$ debugging
      & Mimics human problem-solving cycle; example retrieval
      & 93.9\% & 83.1\% \\
    \bottomrule
  \end{tabular}
\end{table*}

\paragraph{ChatDev~\cite{qian2024chatdev}.}
From Tsinghua University and collaborators, ChatDev models a virtual software
company in which agents assume executive and technical roles: CEO, Chief
Product Officer (CPO), Chief Technology Officer (CTO), Programmer, Reviewer,
Tester, and Art Designer. The system organises work into a waterfall-like
sequence of ``chat chains'': each phase (designing, coding, testing,
documenting) corresponds to a multi-turn conversation among relevant agents.
For example, the CPO and CEO refine requirements, then the CTO and Programmer
generate code, then the Programmer and Tester validate it. ChatDev introduces
\emph{communicative dehallucination}: agents maintain memory of earlier
discussions and replay relevant context to reduce inconsistencies and
hallucinations across phases. The framework also supports role-playing prompts
that instruct each LLM instantiation to adopt a persona, though all agents
share the same base model. ChatDev demonstrates end-to-end software generation
from natural-language descriptions to executable programs, but performance
metrics on standardised benchmarks are limited in the original publication.

\paragraph{MetaGPT~\cite{hong2024metagpt}.}
MetaGPT, from DeepWisdom and collaborators, encodes human-like Standardized
Operating Procedures (SOPs) into LLM prompts to streamline multi-agent
collaboration. It assigns four primary roles: Product Manager, Architect,
Engineer, and QA. Each role produces structured artifacts---Product Requirement
Documents (PRD), system designs, implementation code, and test plans---passed
sequentially along an ``assembly line''. The key innovation is that agents
communicate via formalised documents rather than free-form chat, reducing
ambiguity and cascading hallucinations. MetaGPT includes an \emph{executable
feedback module}: generated code is run against tests, and failures trigger
targeted revisions by the Engineer agent. On HumanEval and MBPP, MetaGPT
achieves 85.9\% and 87.7\% pass@1 respectively, outperforming many single-agent
baselines and prior multi-agent systems. The framework demonstrates that
structured inter-agent protocols can significantly reduce errors compared to
informal chat-based collaboration.

\paragraph{AgentCoder~\cite{huang2023agentcoder}.}
AgentCoder, from the University of Hong Kong and collaborators, simplifies role
structure to three agents: Programmer, Test Designer, and Test Executor. The
Programmer generates code. The Test Designer independently generates test
cases, ensuring tests are not biased by the code. The Test Executor runs code
against tests and reports failures. If tests fail, the executor provides
feedback (error messages, stack traces) to the Programmer, which revises the
code. This cycle iterates until all tests pass or a budget is exhausted.
AgentCoder's strength lies in \emph{independent test generation}: because tests
are written without seeing the code, they better reflect specification
requirements rather than implementation quirks. On HumanEval and MBPP,
AgentCoder with GPT-3.5 reaches 77.4\% and 89.1\% pass@1, surpassing
single-agent GPT-3.5 baselines (69.5\%, 63.0\%) and demonstrating that even a
minimal three-agent decomposition improves reliability.

\paragraph{MapCoder~\cite{islam2024mapcoder}.}
MapCoder, from Bangladesh University of Engineering and UCLA, targets
competitive programming by mimicking the human problem-solving cycle through
four agents: Recall, Plan, Code, and Debug. The Recall agent retrieves similar
past problems or solution templates from a knowledge base. The Plan agent
decomposes the problem into algorithmic steps. The Code agent generates the
implementation. The Debug agent runs the code, identifies failures, and
proposes fixes. Agents operate sequentially, passing intermediate results
forward, and backtracking occurs when debugging reveals a flaw in earlier
stages. MapCoder achieves strong results on competitive benchmarks: 93.9\%
pass@1 on HumanEval, 83.1\% on MBPP, and 28.5\% on CodeContests. The framework
shows that task-specific role design (e.g., retrieval for competitive
programming) can yield state-of-the-art performance when aligned with domain
workflows.

\subsection{When Does Multi-Agency Help?}
\label{sec:multi-when}

Multi-agent designs introduce coordination overhead, increased token costs, and
additional failure modes. Whether they outperform monolithic agents depends on
task complexity, the quality of role decomposition, and the effectiveness of
inter-agent communication.

\paragraph{Benefits: specialisation and modularity.}
The primary advantage of multi-agent systems is \emph{specialisation}. When a
single agent must simultaneously plan, code, test, and debug, its attention is
divided and errors in one phase can cascade unchecked. Decomposing these
concerns into separate agents allows each to focus on a narrower task. For
example, AgentCoder's independent test designer generates tests that are more
comprehensive and less biased than tests written by the same agent that wrote
the code. Similarly, MetaGPT's structured artifacts (PRDs, design documents)
force explicit planning before coding, reducing the likelihood that
implementation drifts from requirements. Modularity also enables targeted
improvement: if test generation is weak, only the test-designer agent need be
upgraded, rather than retraining or reprompting the entire system.

\paragraph{Benefits: error detection and correction.}
Multi-agent systems create multiple checkpoints where errors can be caught. A
reviewer agent can identify code smells or logic bugs before execution. A
separate test executor can detect runtime failures and provide diagnostic
feedback to the programmer without the programmer needing to self-critique.
This separation mirrors human code review and pair programming, which are
empirically known to improve code quality. MapCoder's debug agent, for example,
systematically analyses test failures and proposes fixes, iterating until the
code passes. Single-agent systems can approximate this through
self-reflection~\cite{madaan2023selfrefine}, but multi-agent designs make the
separation explicit and avoid the risk of self-confirming biases.

\paragraph{Costs: coordination and token overhead.}
The main cost is coordination complexity. Each inter-agent communication
requires prompt construction, LLM inference, and state synchronisation. If
agents operate sequentially, total latency is the sum of all agent runtimes. If
agents operate in parallel (e.g., multiple programmers proposing solutions),
orchestration logic must merge or vote on outputs. Token costs scale with the
number of agents and communication rounds: MetaGPT and ChatDev can consume
several times the tokens of a single-agent baseline. In practice, frameworks
mitigate this by caching shared context, compressing intermediate artifacts, or
limiting the number of iterations. However, for simple tasks where a
single-agent solution suffices, the overhead is not justified.

\paragraph{Failure modes: handoff errors and misalignment.}
Multi-agent systems introduce new failure modes. Handoff errors occur when one
agent's output is misunderstood or ignored by the next. For example, if the
architect agent specifies an API contract in natural language and the
programmer misinterprets it, the resulting code may be internally consistent
but wrong. MetaGPT's structured artifacts reduce this risk by enforcing
explicit schemas, but free-form chat-based systems like ChatDev remain
vulnerable. Another failure mode is \emph{agent misalignment}: if the test
designer agent generates trivial or incorrect tests, the system may converge on
code that passes those weak tests but fails in production. AgentCoder addresses
this by ensuring tests are written independently, but the quality still depends
on the test-designer prompt and model.

\paragraph{Diminishing returns with agent count.}
Adding more agents does not monotonically improve performance. Beyond a certain
point, coordination overhead dominates and additional roles contribute little.
Empirical studies on ChatDev and MetaGPT suggest that 3--5 agents (planning,
coding, testing, review) capture most of the benefit, while systems with 10+
agents show marginal improvements at much higher cost. The optimal agent count
depends on task complexity: simple function-writing benefits from
programmer-tester pairs, while full software projects may justify
architect-programmer-tester-reviewer-documenter chains. However, even for
complex tasks, the evidence does not yet support very large agent counts
analogous to human teams of dozens.

\paragraph{When to choose multi-agent designs.}
Multi-agent frameworks are most beneficial when: (i)~the task involves multiple
distinct phases (design, implementation, validation) that benefit from
specialisation; (ii)~intermediate artifacts (specifications, test suites) can
be produced and validated independently; (iii)~the base LLM is strong enough
that role-specific prompting yields measurably different behaviours; and
(iv)~the cost overhead is acceptable given the quality improvement. For
simple one-off code snippets, single-agent systems with self-refinement suffice
and are faster. For end-to-end software engineering tasks on benchmarks like
SWE-bench, the evidence is mixed: some multi-agent systems underperform
well-tuned single-agent scaffolds (recall Agentless outperforming many agentic
designs in \S\ref{sec:e2e}), suggesting that scaffold quality matters as much
as agent count. The frontier remains open, and multi-agent research continues
to refine orchestration protocols and role definitions.

% if needed.

\section{Multi-Agent Frameworks for Code Generation}
\label{sec:multiagent}

An alternative architectural approach to code generation distributes the work
across multiple specialised agents, each responsible for a distinct phase or
concern. Where single-agent systems integrate planning, generation, testing,
and refinement into one LLM conversation, multi-agent frameworks decompose the
software lifecycle into roles---architect, programmer, tester, reviewer---and
orchestrate collaboration among agents playing those roles. This
division-of-labour mirrors human software teams and aims to improve code
quality, modularity, and error detection beyond what monolithic agents achieve.

This section surveys role decomposition and communication patterns
(\S\ref{sec:multi-roles}), examines representative framework designs
(\S\ref{sec:multi-frameworks}), and assesses when multi-agent designs
outperform monolithic alternatives (\S\ref{sec:multi-when}).

\subsection{Role Decomposition and Communication}
\label{sec:multi-roles}

Multi-agent systems partition code-generation tasks along two primary
dimensions: \emph{temporal stages} (design, implementation, testing, review)
and \emph{functional concerns} (requirements analysis, algorithm design,
execution, debugging). Agents are assigned roles corresponding to these
partitions, and collaboration protocols define how information flows among them.

\paragraph{Role taxonomy.}
Common roles mirror human software teams. A \emph{product manager} or
\emph{architect} agent refines high-level requirements into technical
specifications. A \emph{programmer} or \emph{coder} agent generates
implementation. A \emph{test designer} agent writes test cases independently of
the code. A \emph{test executor} agent runs code against tests and reports
failures. A \emph{reviewer} agent critiques code for correctness, style, or
security. More specialised roles include a \emph{retrieval agent} that searches
documentation or past solutions, and a \emph{debugger agent} that localises and
repairs faults. Not all frameworks use all roles; designs range from three
agents to dozens.

\paragraph{Communication patterns.}
Agents exchange information through structured protocols. The simplest is
\emph{sequential handoff}: the architect produces a design document, passes it
to the programmer, who generates code, passes it to the tester, and so on. More
sophisticated patterns include \emph{iterative feedback loops}, where the test
executor returns failure traces to the programmer for revision, and
\emph{debate or voting}, where multiple agents propose alternatives and a
consensus mechanism selects the best. Communication may use natural language
(agents converse as humans would), structured artifacts (JSON schemas, API
specifications, test suites), or hybrid approaches.

\paragraph{Memory and state management.}
Multi-agent systems maintain shared state through blackboard architectures,
message queues, or centralised repositories. A blackboard allows all agents to
read and write a common workspace containing the current design, code, tests,
and execution logs. Message queues enforce ordering and ensure agents see
consistent snapshots of the task. Centralised repositories (e.g., a simulated
Git workflow) provide versioning and rollback. The choice of state management
affects coordination overhead: blackboards are flexible but require
synchronisation; message queues are deterministic but rigid; repositories add
version-control semantics but increase complexity.

\subsection{Framework Designs}
\label{sec:multi-frameworks}

We survey four representative multi-agent frameworks that span diverse design
choices in orchestration, role granularity, and evaluation focus.
Table~\ref{tab:multiagent} provides a comparative overview.

\begin{table*}[t]
  \centering
  \caption{Comparison of multi-agent code-generation frameworks.}
  \label{tab:multiagent}
  \small
  \begin{tabular}{@{}llp{0.28\textwidth}p{0.25\textwidth}cc@{}}
    \toprule
    \textbf{System} & \textbf{Roles} & \textbf{Communication} & \textbf{Key mechanism} & \textbf{HumanEval} & \textbf{MBPP} \\
    \midrule
    ChatDev & CEO, CPO, CTO, Programmer, Reviewer, Tester
      & Sequential chat-based handoffs with communicative dehallucination
      & Waterfall-inspired phases; memory replay to prevent hallucinations
      & -- & -- \\
    MetaGPT & Product manager, Architect, Engineer, QA
      & SOP-encoded prompts with structured artifacts (PRD, design docs)
      & Assembly-line paradigm; executable feedback module
      & 85.9\% & 87.7\% \\
    AgentCoder & Programmer, Test designer, Test executor
      & Iterative loop: code $\to$ tests $\to$ execution $\to$ feedback
      & Independent test generation; execution-driven refinement
      & 77.4\% & 89.1\% \\
    MapCoder & Retrieval, Planner, Coder, Debugger
      & Sequential pipeline with recalling $\to$ planning $\to$ coding $\to$ debugging
      & Mimics human problem-solving cycle; example retrieval
      & 93.9\% & 83.1\% \\
    \bottomrule
  \end{tabular}
\end{table*}

\paragraph{ChatDev~\cite{qian2024chatdev}.}
From Tsinghua University and collaborators, ChatDev models a virtual software
company in which agents assume executive and technical roles: CEO, Chief
Product Officer (CPO), Chief Technology Officer (CTO), Programmer, Reviewer,
Tester, and Art Designer. The system organises work into a waterfall-like
sequence of ``chat chains'': each phase (designing, coding, testing,
documenting) corresponds to a multi-turn conversation among relevant agents.
For example, the CPO and CEO refine requirements, then the CTO and Programmer
generate code, then the Programmer and Tester validate it. ChatDev introduces
\emph{communicative dehallucination}: agents maintain memory of earlier
discussions and replay relevant context to reduce inconsistencies and
hallucinations across phases. The framework also supports role-playing prompts
that instruct each LLM instantiation to adopt a persona, though all agents
share the same base model. ChatDev demonstrates end-to-end software generation
from natural-language descriptions to executable programs, but performance
metrics on standardised benchmarks are limited in the original publication.

\paragraph{MetaGPT~\cite{hong2024metagpt}.}
MetaGPT, from DeepWisdom and collaborators, encodes human-like Standardized
Operating Procedures (SOPs) into LLM prompts to streamline multi-agent
collaboration. It assigns four primary roles: Product Manager, Architect,
Engineer, and QA. Each role produces structured artifacts---Product Requirement
Documents (PRD), system designs, implementation code, and test plans---passed
sequentially along an ``assembly line''. The key innovation is that agents
communicate via formalised documents rather than free-form chat, reducing
ambiguity and cascading hallucinations. MetaGPT includes an \emph{executable
feedback module}: generated code is run against tests, and failures trigger
targeted revisions by the Engineer agent. On HumanEval and MBPP, MetaGPT
achieves 85.9\% and 87.7\% pass@1 respectively, outperforming many single-agent
baselines and prior multi-agent systems. The framework demonstrates that
structured inter-agent protocols can significantly reduce errors compared to
informal chat-based collaboration.

\paragraph{AgentCoder~\cite{huang2023agentcoder}.}
AgentCoder, from the University of Hong Kong and collaborators, simplifies role
structure to three agents: Programmer, Test Designer, and Test Executor. The
Programmer generates code. The Test Designer independently generates test
cases, ensuring tests are not biased by the code. The Test Executor runs code
against tests and reports failures. If tests fail, the executor provides
feedback (error messages, stack traces) to the Programmer, which revises the
code. This cycle iterates until all tests pass or a budget is exhausted.
AgentCoder's strength lies in \emph{independent test generation}: because tests
are written without seeing the code, they better reflect specification
requirements rather than implementation quirks. On HumanEval and MBPP,
AgentCoder with GPT-3.5 reaches 77.4\% and 89.1\% pass@1, surpassing
single-agent GPT-3.5 baselines (69.5\%, 63.0\%) and demonstrating that even a
minimal three-agent decomposition improves reliability.

\paragraph{MapCoder~\cite{islam2024mapcoder}.}
MapCoder, from Bangladesh University of Engineering and UCLA, targets
competitive programming by mimicking the human problem-solving cycle through
four agents: Recall, Plan, Code, and Debug. The Recall agent retrieves similar
past problems or solution templates from a knowledge base. The Plan agent
decomposes the problem into algorithmic steps. The Code agent generates the
implementation. The Debug agent runs the code, identifies failures, and
proposes fixes. Agents operate sequentially, passing intermediate results
forward, and backtracking occurs when debugging reveals a flaw in earlier
stages. MapCoder achieves strong results on competitive benchmarks: 93.9\%
pass@1 on HumanEval, 83.1\% on MBPP, and 28.5\% on CodeContests. The framework
shows that task-specific role design (e.g., retrieval for competitive
programming) can yield state-of-the-art performance when aligned with domain
workflows.

\subsection{When Does Multi-Agency Help?}
\label{sec:multi-when}

Multi-agent designs introduce coordination overhead, increased token costs, and
additional failure modes. Whether they outperform monolithic agents depends on
task complexity, the quality of role decomposition, and the effectiveness of
inter-agent communication.

\paragraph{Benefits: specialisation and modularity.}
The primary advantage of multi-agent systems is \emph{specialisation}. When a
single agent must simultaneously plan, code, test, and debug, its attention is
divided and errors in one phase can cascade unchecked. Decomposing these
concerns into separate agents allows each to focus on a narrower task. For
example, AgentCoder's independent test designer generates tests that are more
comprehensive and less biased than tests written by the same agent that wrote
the code. Similarly, MetaGPT's structured artifacts (PRDs, design documents)
force explicit planning before coding, reducing the likelihood that
implementation drifts from requirements. Modularity also enables targeted
improvement: if test generation is weak, only the test-designer agent need be
upgraded, rather than retraining or reprompting the entire system.

\paragraph{Benefits: error detection and correction.}
Multi-agent systems create multiple checkpoints where errors can be caught. A
reviewer agent can identify code smells or logic bugs before execution. A
separate test executor can detect runtime failures and provide diagnostic
feedback to the programmer without the programmer needing to self-critique.
This separation mirrors human code review and pair programming, which are
empirically known to improve code quality. MapCoder's debug agent, for example,
systematically analyses test failures and proposes fixes, iterating until the
code passes. Single-agent systems can approximate this through
self-reflection~\cite{madaan2023selfrefine}, but multi-agent designs make the
separation explicit and avoid the risk of self-confirming biases.

\paragraph{Costs: coordination and token overhead.}
The main cost is coordination complexity. Each inter-agent communication
requires prompt construction, LLM inference, and state synchronisation. If
agents operate sequentially, total latency is the sum of all agent runtimes. If
agents operate in parallel (e.g., multiple programmers proposing solutions),
orchestration logic must merge or vote on outputs. Token costs scale with the
number of agents and communication rounds: MetaGPT and ChatDev can consume
several times the tokens of a single-agent baseline. In practice, frameworks
mitigate this by caching shared context, compressing intermediate artifacts, or
limiting the number of iterations. However, for simple tasks where a
single-agent solution suffices, the overhead is not justified.

\paragraph{Failure modes: handoff errors and misalignment.}
Multi-agent systems introduce new failure modes. Handoff errors occur when one
agent's output is misunderstood or ignored by the next. For example, if the
architect agent specifies an API contract in natural language and the
programmer misinterprets it, the resulting code may be internally consistent
but wrong. MetaGPT's structured artifacts reduce this risk by enforcing
explicit schemas, but free-form chat-based systems like ChatDev remain
vulnerable. Another failure mode is \emph{agent misalignment}: if the test
designer agent generates trivial or incorrect tests, the system may converge on
code that passes those weak tests but fails in production. AgentCoder addresses
this by ensuring tests are written independently, but the quality still depends
on the test-designer prompt and model.

\paragraph{Diminishing returns with agent count.}
Adding more agents does not monotonically improve performance. Beyond a certain
point, coordination overhead dominates and additional roles contribute little.
Empirical studies on ChatDev and MetaGPT suggest that 3--5 agents (planning,
coding, testing, review) capture most of the benefit, while systems with 10+
agents show marginal improvements at much higher cost. The optimal agent count
depends on task complexity: simple function-writing benefits from
programmer-tester pairs, while full software projects may justify
architect-programmer-tester-reviewer-documenter chains. However, even for
complex tasks, the evidence does not yet support very large agent counts
analogous to human teams of dozens.

\paragraph{When to choose multi-agent designs.}
Multi-agent frameworks are most beneficial when: (i)~the task involves multiple
distinct phases (design, implementation, validation) that benefit from
specialisation; (ii)~intermediate artifacts (specifications, test suites) can
be produced and validated independently; (iii)~the base LLM is strong enough
that role-specific prompting yields measurably different behaviours; and
(iv)~the cost overhead is acceptable given the quality improvement. For
simple one-off code snippets, single-agent systems with self-refinement suffice
and are faster. For end-to-end software engineering tasks on benchmarks like
SWE-bench, the evidence is mixed: some multi-agent systems underperform
well-tuned single-agent scaffolds (recall Agentless outperforming many agentic
designs in \S\ref{sec:e2e}), suggesting that scaffold quality matters as much
as agent count. The frontier remains open, and multi-agent research continues
to refine orchestration protocols and role definitions.
